\chapter[Jesus nimmt die Sünder an!]{„Jesus nimmt die Sünder an!“}
\quellseite{5}

Die köstliche Schriftwahrheit, daß Jesus die Sünder annimmt, Luc. 15, 2., dient nicht nur dazu, uns Christen den einzig wahren Trost in Sündennoth zu bringen, sondern sie treibt uns auch an, das Evangelium von Christo, dem Sünderheiland, andern Sündern zu bringen. Je mehr und besser wir diesen Trost in unserm eigenen Herzen empfinden, desto freudiger und eifriger werden wir, ihn andern zu bringen.

Ich will dir, lieber Leser, kurz von zwei Sündern erzählen, die unser lieber Herr Jesus noch in der letzten Stunde angenommen hat. Du wirst dich dann von Herzen freuen und die Mission an solchen Sündern durch deine Fürbitte unterstützen.

Bei meinem wöchentlichen Besuch im Zuchthause zu Chester, Ill., traf ich im Hospital einen jungen Gefangenen, einen Mörder, der zu lebenslänglicher Haft verurtheilt war. Er lag an der Schwindsucht danieder und konnte nicht mehr lange leben. Niemand kümmerte sich um sein Seelenheil.

Ich begrüßte ihn und sagte zu ihm: „G., du bist schwer krank und wirst bald sterben.“

Er: „Ja, Herr Pastor.“

Ich: „Wohin wirst du denn fahren, wenn du stirbst?“

Er: „O, ich fahre zur Hölle, das ist gewiß.“

Ich: „Woher weißt du denn, daß du zur Hölle fährst?“

Er: „Ich bin ein Mörder, und Gott muß mich verdammen. Ich habe schrecklich gesündigt und habe nichts anderes zu erwarten.“

Ich: „Gibt es denn nichts, das dich retten könnte? Hast du noch nicht vom Heiland der Sünder gehört?“

Er: „Ich bin getauft, aber nicht christlich erzogen und habe nie gehört, daß es einen Heiland der Sünder gibt.“

Ich: „Soll ich dir vom Heiland erzählen?“

Er: „O bitte, Herr Pastor.“

Ich erzählte diesem armen Menschen nun von Jesu, von seiner Geburt und seinem Erlösungswerk. Ich sagte ihm, daß Jesus alle Sünder vollkommen erlöst habe und daß sie durch sein Verdienst könnten und sollten ewig selig sein, wenn sie nur dieses glaubten. Ich erzählte ihm auch vom Schächer am Kreuz, der noch in der letzten Stunde sich zu Jesu gewendet und sich im Glauben fest auf Jesu Erlösung verlassen und darauf das Verlangen, in das Paradies einzugehen, gegründet habe. Das verklärte Gesicht, der hoffnungsvolle Blick des armen Menschen läßt sich nicht beschreiben. Er bat mich dringend, ihn bald wieder zu besuchen.

Etwa fünf Wochen lang besuchte ich ihn jede Woche. Seine Begierde, das Wort Gottes zu hören, war groß und seine Freude über den darin enthaltenen Trost noch größer. Immer wieder mußte ich ihm aus der Schrift vorlesen.

Die letzten beiden Male mußte ich ihn im Sterbezimmer, in welches er gebracht worden war, besuchen. Er wußte aber nichts mehr von Furcht vor dem Tode. Er bekannte seine Sünden, aber er that es mit der fröhlichen Gewißheit, daß Jesus sie alle gebüßt habe. Beim letzten Abschied sagte er: „Herr Pastor, ich will gerne sterben, ich gehe dann heim.“ Als ich wiederkam, war er heimgegangen, und ich konnte mir sagen, daß der liebe Herr Jesus den armen Sünder angenommen habe.

Doch Jesus wollte durch meine geringen Dienste noch einen andern Sünder retten. Davon laß mich dir noch kurz erzählen, lieber Leser.

Im Sterbezimmer neben G. lag ein Neger, auch ein Mörder, ebenfalls im letzten Stadium der Schwindsucht. Als G. gestorben war, sagte ich zu dem Neger: „Sam, G. ist gestorben, wer wird nun wohl der nächste sein?“ Mühsam richtete er sich auf und antwortete: „Ich, Herr Pastor.“

Ich: „Und wie wirst du fahren?“

Er: „O Herr Pastor, denken Sie ja nicht, daß ich nicht gehört habe, was Sie zu G. gesagt haben: ‚Das Blut Jesu Christi, des Sohnes Gottes, macht uns rein von aller Sünde.‘ Ich bin ein schrecklicher Sünder, einen größeren gibt es nicht; aber das Blut Jesu Christi macht mich rein von allen Sünden.“

Ich: „Sam, wie weißt du denn, daß das wahr ist?“

Er: „Das kleine Buch“ (das Neue Testament), „aus dem Sie vorgelesen haben, ist Gottes Wort. Gott sagt es und Gott lügt nicht.“

Ich: „Wie weißt du denn, daß du damit gemeint bist?“

Er: „Haben Sie nicht gelesen: ‚Also hat Gott die Welt geliebet, daß er seinen eingebornen Sohn gab, auf daß alle, die an ihn glauben, nicht verloren werden, sondern das ewige Leben haben‘? Gehört Sam nicht zur Welt? Ist er nicht einer der Schlimmsten in der Welt? Gewiß bin ich gemeint; das Blut Jesu Christi, des Sohnes Gottes, macht mich rein.“

Ich: „Sam, wenn die Todesstunde kommt, wird der Teufel versuchen, dir solchen Glauben zu rauben.“

Er: „O Herr Pastor, er kann nicht, denn Jesus hat versprochen, bei mir zu sein bis ans Ende, und Jesus ist stärker als der Teufel.“

Ich ermunterte ihn dann und tröstete ihn, daß Gott ihn bei diesem Glauben erhalten und ihn selig machen werde. Zum Abschied sagte er zu mir: „Herr Pastor, ich gehe heim, bald gehe ich heim, denn das Blut Jesu Christi, des Sohnes Gottes, macht mich rein von allen Sünden. Wenn Sie einmal auch in den Himmel kommen, dann will ich Sie zum Heiland bringen und sagen: Dies ist der deutsche Pastor, von dem habe ich es gehört: ‚Das Blut Jesu Christi‘ \&c., und dadurch bin ich selig geworden.
\quellseite{6}
Dann wollen wir uns aber freuen!“
Er bat mich dann noch, ihn wieder zu besuchen, wenn er noch lebe, setzte aber gleich hinzu, er werde schon heimgegangen sein, wenn ich wiederkommen würde. Als ich wiederkam, war Sam entschlafen. Sein Wächter berichtete mir, daß er an demselben Abend gegen zehn Uhr fast eine ganze Stunde laut gebetet und mit Gott gerungen habe, worauf er sich zurückgelegt habe und ruhig eingeschlafen sei.

Sieh, lieber Leser: „Jesus nimmt die Sünder an!“ Auch die allergrößten Sünder nimmt er an, oft noch in der allerletzten Stunde. Ist es nicht köstlich, dem Herrn Jesu dienen zu dürfen, den Sündern, die er annehmen will, das rettende Evangelium zu bringen? Willst du nicht helfen? O schließe die Mission, welche an den Gefangenen betrieben wird, ein in deine Fürbitte. Später kann ich dir einmal wieder erzählen von der großen Gnade Gottes, die er in dieser Mission erweist. E. Strn.
